% 原章分章页32—36，第2.5.2节。
\subsection{线性与非线性混合结构方程模型}

前面讨论了线性结构方程模型和非线性结构方程模型。本节介绍具有高斯噪声的线性结构方程模型，或线性与非线性混合结构方程模型（Ernest et al. 2016）。一般而言，网络中某些边的函数可能是线性的，另一些边的函数则可能是非线性的。线性与非线性混合结构方程模型是最一般的因果模型。例如，图2.7给出了一个混合模型，其中实线表示非线性边，虚线表示线性边。
\begin{figure}[htbp]
\centering\begin{tikzpicture}[x=1.4cm,y=1cm,>=Stealth,
ch02v/.style={ellipse,draw=inkblue,minimum width=9mm,minimum height=7mm,font=\small},
ch02non/.style={->,red,semithick},
ch02lin/.style={->,green!60!black,dashed,semithick}]
\node[ch02v](ch0207n1)at(-1,3.5){1};
\node[ch02v](ch0207n2)at(1,3.5){2};
\node[ch02v](ch0207n3)at(3,3.5){3};
\node[ch02v](ch0207n4)at(0,2){4};
\node[ch02v](ch0207n5)at(2.7,2){5};
\node[ch02v](ch0207n6)at(1.5,.5){6};
\node[ch02v](ch0207n7)at(1.5,-1){7};
\draw[ch02non](ch0207n1)--(ch0207n4);
\draw[ch02lin](ch0207n2)--(ch0207n4);
\draw[ch02lin](ch0207n3)--(ch0207n2);
\draw[ch02non](ch0207n2)--(ch0207n5);
\draw[ch02lin](ch0207n3)--(ch0207n5);
\draw[ch02lin](ch0207n4)--(ch0207n6);
\draw[ch02non](ch0207n5)--(ch0207n6);
\draw[ch02non](ch0207n6)--(ch0207n7);
\node[font=\small,anchor=east](ch0207txt1)at(-1.5,2.6){非线性边};
\draw[->,blue!65](ch0207txt1.east)--(-.55,2.9);
\node[font=\small,anchor=west](ch0207txt2)at(3.7,2.7){线性边};
\draw[->,blue!65](ch0207txt2.west)--(2.9,2.8);
\end{tikzpicture}

\origfig{7}{线性与非线性混合结构方程模型示意图。}
\end{figure}
式（2.77）中的一般非线性结构方程模型可以写为（Ernest et al. 2016）
\begin{equation}
Y_j=\mu_j+\sum_{i\in\pa_D(j)}f_{ji}(W_i)+\varepsilon_j,
\tag{2.92}
\end{equation}
其中，$\mu_j\in\R$，$f_{ji}$ 二次可微，$\E[f_{ji}(W_i)]=0$，$\varepsilon_j\sim N(0,\sigma_j^2)$，$\sigma_j^2>0$，$j=1,2,\ldots,p$。图 $D$ 中有向边 $i\to j$ 对应的函数 $f_{ji}$ 可以是线性函数，也可以是非线性函数。如果 $f_{ji}$ 是线性函数，就称边 $i\to j$ 为线性边；否则称为非线性边。正如 Ernest et al.（2016）所指出的，边的线性或非线性是相对于某个具体 DAG $D$ 定义的。

总因果效应可以分解为线性因果效应与非线性因果效应。模型（2.92）可以改写为
\begin{equation}
y_j=\mu_j+\sum_{l=1}^{l_j}\alpha_{jl}w_l
+\sum_{i=1}^{i_j}f_{ji}(w_i)+\varepsilon_j,
\tag{2.93}
\end{equation}
其中，$l_j$ 表示节点 $j$ 通过线性边连接的父节点个数，$i_j$ 表示节点 $j$ 通过非线性边连接的父节点个数。例如，节点4的函数模型为
\[
y_4=\mu_4+\alpha_{41}y_1+f_{42}(y_2)+\varepsilon_4.
\]
一般而言，函数 $f_{ji}$ 未知，需要用非参数函数逼近。本节仍使用平滑样条表示非线性函数 $f_{ji}$。考虑平滑样条回归模型（附录2.A）：
\begin{equation}
y_j=\mu_j+\sum_{l=1}^{l_j}\alpha_{jl}w_l
+\sum_{i=1}^{i_j}L_{ji}f_{ji}+\varepsilon_j,
\tag{2.94}
\end{equation}
其中，$L_{ji}$ 是定义在 RKHS $H^i$ 上的连续泛函，$f_{ji}$ 表示一般未知函数，$\varepsilon_i$ 为均值为零、共同方差为 $\sigma_j^2$ 的独立随机误差。希尔伯特空间 $H^i$ 可分解为两个子空间（Wang 2011）：
\[
H^i=H_0^i\oplus H_1^i,
\]
其中，$H_0^i$ 为有限维空间，具有正交归一基函数 $\varphi_1^i(x),\ldots,\varphi_p^i(x)$；它的正交补 $H_1^i$ 是具有再生核 $R_1^i(x,z)$ 的 RKHS。

估计式（2.94）中的回归系数 $\alpha_{jl}$ 和函数 $f_{ji}$ 时，应使观测值 $Y_j$ 与估计量 $\widehat\alpha_{ji},\widehat f_{ji}$ 之间的误差，以及对光滑性的偏离，都尽可能小。用函数 $f_{ji}$ 向 RKHS $H_1$ 的投影 $P_1f_{ji}$ 作为惩罚函数。为此，回归系数估计量和函数 $f_{ji}$ 的平滑样条估计量所对应的目标函数为
\begin{equation}
\frac1n\sum_{v=1}^{n}
\left(y_j^v-\sum_{l=1}^{l_j}\alpha_{jl}w_l^v
-\sum_{i=1}^{i_j}L_vf_{ji}\right)^2
+\lambda\left\lVert\sum_{i=1}^{i_j}P_1f_{ji}\right\rVert^2,
\tag{2.95}
\end{equation}
其中，$\lambda$ 为平滑惩罚参数，用于平衡拟合优度和光滑性。令 $R^i(w,z)$ 为与 $f_{ji}$ 对应的再生核，可分解为
\[
R^i(w,z)=R_0^i(w,z)+R_1^i(w,z),
\]
其中
\begin{equation}
R_0^i(w,z)=\sum_{m=1}^{p}\varphi_m^i(w)\varphi_m^i(z).
\tag{2.96}
\end{equation}
假设估计量 $\widehat f_{ji}$ 表示为
\begin{equation}
\widehat f_{ji}(w)=\sum_{m=1}^{p}a_m^i\varphi_m^i(w)
+\sum_{u=1}^{n}b_u^i\xi_u^i(w),
\tag{2.97}
\end{equation}
其中 $\xi_u^i(w)=L_{u(z)}R_1^i(w,z)$。因此，
\begin{equation}
L_{v(w)}\widehat f_{ji}(w)=
\sum_{m=1}^{p}a_m^iL_{v(w)}\varphi_m^i(w)
+\sum_{u=1}^{n}b_u^iL_{v(w)}L_{u(z)}R_1^i(w,z),
\tag{2.98}
\end{equation}
或者
\[
\begin{aligned}
\begin{bmatrix}
L_{1(w)}\widehat f_{ji}(w)\\\vdots\\L_{n(x)}\widehat f(x)
\end{bmatrix}
={}&
\begin{bmatrix}
L_{1(w)}\varphi_1^i(w)&\cdots&L_{1(w)}\varphi_p^i(w)\\
\vdots&\vdots&\vdots\\
L_{n(w)}\varphi_1^i(w)&\cdots&L_{n(w)}\varphi_p^i(w)
\end{bmatrix}
\begin{bmatrix}a_1^i\\\vdots\\a_p^i\end{bmatrix}\\
&{}+
\begin{bmatrix}
L_{1(w)}L_{1(z)}R_1^i(x,z)&\cdots&L_{1(w)}L_{n(z)}R_1^i(x,z)\\
\vdots&\vdots&\vdots\\
L_{n(w)}L_{1(z)}R_1^i(x,z)&\cdots&L_{n(w)}L_{n(z)}R_1^i(x,z)
\end{bmatrix}
\begin{bmatrix}b_1^i\\\vdots\\b_n^i\end{bmatrix}\\
={}&T^ia^i+\Sigma^ib^i,
\end{aligned}
\]
其中
\[
T^i=
\begin{bmatrix}
L_{1(w)}\varphi_1^i(w)&\cdots&L_{1(w)}\varphi_p^i(w)\\
\vdots&\vdots&\vdots\\
L_{n(w)}\varphi_1^i(w)&\cdots&L_{n(w)}\varphi_p^i(w)
\end{bmatrix},
\qquad
a^i=\begin{bmatrix}a_1^i\\\vdots\\a_p^i\end{bmatrix},
\]
\[
\Sigma^i=
\begin{bmatrix}
L_{1(w)}L_{1(z)}R_1^i(x,z)&\cdots&L_{1(w)}L_{n(z)}R_1^i(x,z)\\
\vdots&\vdots&\vdots\\
L_{n(w)}L_{1(z)}R_1^i(x,z)&\cdots&L_{n(w)}L_{n(z)}R_1^i(x,z)
\end{bmatrix},
\qquad
b^i=\begin{bmatrix}b_1^i\\\vdots\\b_n^i\end{bmatrix}.
\]
令
\[
\begin{gathered}
Y_j=\begin{bmatrix}y_j^1\\\vdots\\y_j^n\end{bmatrix},
\quad
W=\begin{bmatrix}w_1^1&\cdots&w_{l_j}^1\\
\vdots&\vdots&\vdots\\w_1^n&\cdots&w_{l_j}^n\end{bmatrix},
\quad
\alpha^j=\begin{bmatrix}\alpha_{j1}\\\vdots\\\alpha_{jl_j}\end{bmatrix},\\
a=\begin{bmatrix}a^1\\\vdots\\a^{i_j}\end{bmatrix},
\quad b=\begin{bmatrix}b^1\\\vdots\\b^{i_j}\end{bmatrix},
\quad\gamma=\begin{bmatrix}\alpha\\a\end{bmatrix},\\
T=[T^1\ \cdots\ T^{i_j}],\qquad A=[W\ T],
\qquad\Sigma=[\Sigma^1\ \cdots\ \Sigma^{i_j}].
\end{gathered}
\]
还可以证明
\begin{equation}
\left|\sum_{i=1}^{i_j}P_1f_{ji}\right|^2=b^T\Sigma b.
\tag{2.99}
\end{equation}
用矩阵形式，式（2.95）可以简化为
\begin{equation}
\min_{\gamma,b}\quad
\frac1n\lVert Y_j-A\gamma-\Sigma b\rVert^2+\lambda b^T\Sigma b.
\tag{2.100}
\end{equation}
优化问题（2.100）的解为（附录2.A）
\begin{equation}
\begin{aligned}
A\gamma+(\Sigma+n\lambda I)b&=Y_j,\\
A^Tb&=0.
\end{aligned}\tag{2.101}
\end{equation}
对矩阵 $A$ 作 QR 分解：
\[
T=[Q_1\ Q_2]\begin{bmatrix}R\\0\end{bmatrix}.
\]
最终得到平滑样条回归系数
\begin{equation}
\begin{aligned}
\widehat\gamma&=
R^{-1}Q_1^T[I-MQ_2(Q_2^TMQ_2)^{-1}Q_2^T]Y_j,\\
\widehat b&=Q_2(Q_2^TMQ_2)^{-1}Q_2^TY_j,
\end{aligned}\tag{2.102}
\end{equation}
其中 $M=\Sigma+n\lambda I$。

现在定义父节点集合同时通过线性边和非线性边连接时的节点评分。假设节点 $j$ 有 $m_j$ 个父节点。考虑有向边 $i\to j$，定义线性边指示变量：
\begin{equation}
\pi_i=
\begin{cases}
1,&\text{边 }i\to j\text{ 为线性边},\\
0,&\text{边 }i\to j\text{ 为非线性边}.
\end{cases}
\tag{2.103}
\end{equation}
函数模型（2.93）可以改写为
\begin{equation}
y_j=\mu_j+\sum_{li=1}^{m_j}\pi_i\alpha_{jl}w_l
+\sum_{i=1}^{im_j}(1-\pi_i)f_{ji}(w_i)+\varepsilon_j.
\tag{2.104}
\end{equation}
为了使模型（2.104）拟合数据，需要极小化以下惩罚目标函数：
\begin{equation}
\begin{aligned}
&\frac1n\sum_{v=1}^{n}
\left(y_j^v-\sum_{i=1}^{m_j}\pi_i\alpha_{jl}w_l^v
-\sum_{i=1}^{m_j}(1-\pi_i)L_vf_{ji}\right)^2\\
&\qquad+\lambda\left\lVert\sum_{i=1}^{m_j}(1-\pi_i)P_1f_{ji}\right\rVert^2.
\end{aligned}\tag{2.105}
\end{equation}
展开连续泛函 $L_v$，式（2.105）可以进一步化为
\begin{equation}
\begin{aligned}
&\frac1n\sum_{v=1}^{n}\Biggl\{
y_j^v-\sum_{i=1}^{m_j}\pi_i\alpha_{ji}w_i^v
-\sum_{i=1}^{m_j}(1-\pi_i)\\
&\quad{}\times\left[
\sum_{q=1}^{p}a_q^iL_{v(w)}\varphi_q^i(w)
+\sum_{u=1}^{n}b_u^iL_{v(w)}L_{v(z)}R_1^i(w,z)
\right]\Biggr\}^2\\
&\qquad+\lambda\left\lVert\sum_{i=1}^{m_j}(1-\pi_i)P_1f_{ji}\right\rVert^2.
\end{aligned}\tag{2.106}
\end{equation}
用矩阵形式，优化问题（2.106）可以改写为
\begin{equation}
\begin{aligned}
\min_{\pi,\alpha,a,b}\quad&
\frac1n\lVert Y_j-W\pi\alpha-T(I-\pi)a-\Sigma(I-\pi)b\rVert^2\\
&{}+\lambda b^T(I-\pi)\Sigma(I-\pi)b,
\end{aligned}\tag{2.107}
\end{equation}
其中 $\pi=\diag(\pi_1,\ldots,\pi_{m_j})$，$I$ 为单位矩阵，其余量的定义同前。在优化问题（2.107）中，对角矩阵 $\pi$ 的元素取 0 或 1，向量 $\alpha,a,b$ 取实数。因此，问题（2.107）是混合整数规划问题，其求解将在第2.6节讨论。求解后，节点 $j$ 的评分函数定义为
\begin{equation}
\begin{aligned}
\operatorname{Score}(j,\pa_j)
={}&\frac1n\lVert Y_j-W\pi\alpha-T(I-\pi)a-\Sigma(I-\pi)b\rVert^2\\
&{}+\lambda b^T(I-\pi)\Sigma(I-\pi)b.
\end{aligned}\tag{2.108}
\end{equation}
